\begin{definitionbox}{Definition (Bounded Subset of the Real Line)}
\begin{definition}[Bounded Subset of the Real Line]
\label{def:bounded-subset-of-r}
A set $A\subseteq\mathbb{R}$ is \emph{bounded} if there exists $M>0$ such that
\[
  (\forall x\in A)(|x|\leq M).
\]
Equivalently,
\[
  A\text{ is bounded}
  \Longleftrightarrow
  (\exists M>0)(\forall x\in A)(-M\leq x\leq M).
\]
\end{definition}
\end{definitionbox}
\begin{remark*}[Standard quantified statement]
\[
E \text{ is bounded} \Longleftrightarrow \exists M\ge 0\;\forall x\in E\;(|x|\le M).
\]
\end{remark*}
\begin{remark*}[Predicate reading]
\[
E\subseteq\mathbb{R},\qquad P=\mathsf{OrderedSet}(\mathbb{R},\leq),\qquad
\operatorname{Bounded}(E,P)
\Longleftrightarrow
\exists M\ge 0\;\forall x\in E\;(|x|\le M).
\]
\end{remark*}
\begin{remark*}[Negated quantified statement]
\[
E \text{ is unbounded} \Longleftrightarrow \forall M\ge 0\;\exists x\in E\;(|x|>M).
\]
\end{remark*}

\begin{remark*}[Negation predicate reading]
\[
E \text{ is unbounded} \Longleftrightarrow \forall M\ge 0\;\exists x\in E\;(|x|>M).
\]

The negation predicate reading is the predicate-level form of the displayed negated statement.
\end{remark*}

\begin{remark*}[Failure modes]
\begin{description}
  \item[Exposition.]
  Failure is witnessed by the displayed negated or contrapositive condition.

  \item[\operatorname{Bounded}.]
  \[
E \text{ is unbounded} \Longleftrightarrow \forall M\ge 0\;\exists x\in E\;(|x|>M).
\]
\[
E\subseteq\mathbb{R},\qquad P=\mathsf{OrderedSet}(\mathbb{R},\leq),\qquad
\operatorname{Bounded}(E,P)
\Longleftrightarrow
\exists M\ge 0\;\forall x\in E\;(|x|\le M).
\]
\end{description}
\end{remark*}
\begin{remark*}[Interpretation]
Boundedness says $E$ fits inside a single closed interval --- it does not run off to infinity. Together with closedness it is one half of the Heine--Borel characterization of compact subsets of $\mathbb{R}$.
\end{remark*}
\begin{dependencies}
\begin{itemize}
  \item \hyperref[def:closed-interval]{Closed Interval}
  \item \hyperref[thm:triangle-inequality]{Triangle Inequality}
\end{itemize}
\end{dependencies}

\subsection*{Set Operations on Intervals}
\label{subsec:set-operations-on-intervals}
